\chapter{User Interface Design}

This chapter describes the user interface design of SHEO, following a standard
user-centred design process: identify the \emph{context of use} $\rightarrow$ derive
\emph{UI requirements} $\rightarrow$ produce a \emph{design} (information architecture
and screen layouts) $\rightarrow$ build a \emph{prototype} $\rightarrow$ conduct
\emph{usability evaluation}. Each section below corresponds to one step and maps its
findings to the system's screens, navigation structures, and interaction patterns.

% -----------------------------------------------------------------------
\section{UI Design Process}
\label{sec:ui-process}

The interface design process followed in SHEO is a five-step user-centred model:

\begin{enumerate}
  \item \textbf{Identify context of use.}  Record who the users are, what environment
        they work in, and on what platform — before any screen is drawn (Section~\ref{sec:ui-context}).
  \item \textbf{Derive UI requirements.}  Translate the context, user goals, and the
        seven golden rules into measurable usability targets and interaction constraints
        (Section~\ref{sec:ui-requirements}).
  \item \textbf{Design.}  Produce an information architecture (sitemap and user flow)
        and low-fidelity wireframes for each screen
        (Sections~\ref{sec:ui-ia} and~\ref{sec:ui-screens}).
  \item \textbf{Prototype.}  Realise the wireframes as a working React single-page
        application; the live implementation is the executable prototype.
  \item \textbf{Evaluate.}  Assess the prototype against the golden rules (heuristic
        walkthrough) and against measurable task criteria (task-based test), then feed
        findings back into the design (Section~\ref{sec:ui-evaluation}).
\end{enumerate}

The process is intentionally iterative: the heuristic pass (step~5 applied to the
initial design) surfaced two concrete improvements that were folded back into
step~3 — surfacing the VND saving estimate \emph{before} the Save button is enabled,
and adding a safety pill on safety-sensitive controls.

% -----------------------------------------------------------------------
\section{Context of Use and User Profiles}
\label{sec:ui-context}

\subsection{Context of use}

SHEO is deployed for a single apartment building, whose units are each occupied by one
resident.  Residents interact with it through a personal computer or a mobile-phone
browser during ordinary domestic hours rather than during a fixed office shift, while the
building owner uses it occasionally to oversee the building.  The platform is a
\textbf{browser-hosted single-page application} (SPA): there is no native install, and
interaction is by mouse and keyboard on desktop or by touch on mobile.  The information
handled during a typical session includes live device power and energy figures, estimated
and measured savings in Vietnamese dong (VND), automation rules, AI-derived
recommendations, and the building-wide electricity tariff.

The project separates \emph{business
stakeholders} who commission and fund the build from \emph{system actors} who operate
the running software (consistent with the stakeholder model of Chapter~1).  The Client (IoT contractor) falls in the first category and does
not log in.  The Customer (building owner or landlord) operates SHEO as the
Administrator actor, who owns no unit of their own; the tenants of the individual units
operate it as Resident actors.

\subsection{User profiles}

Table~\ref{tab:user-profiles} summarises the three operative user classes, their
characteristics, their primary interaction goals, and the consequence of an error for
each class.  The consequence column drives the protection mechanisms described in
Section~\ref{sec:ui-requirements}.

\begin{table}[htbp]
  \caption{User profiles for SHEO.}
  \label{tab:user-profiles}
  \begin{tabularx}{\linewidth}{@{} L{2.2cm} Y Y Y @{}}
    \toprule
    \textbf{User class} & \textbf{Characteristics} & \textbf{Goals / interaction intent}
      & \textbf{Consequence of an error} \\
    \midrule
    Administrator (Customer) &
      Building owner or landlord; owns no unit of their own; non-technical; occasional
      user; motivated by the building's electricity bill. &
      Oversee all units through a read-only building overview; set the building-wide
      tariff and billing cycle; sell, onboard, and offboard units. Never operates devices and
      authors no rules. &
      \textbf{High} — sets the building-wide tariff and onboards/offboards units.
      Safeguarded by explicit confirmations and by read-only oversight of units. \\
    \addlinespace
    Resident (end-user / tenant) &
      Tenant of one unit; mixed ages and abilities; views and controls only their own
      unit's operable devices; sets up personal automation rules. &
      Check own unit's consumption and savings; operate own unit's devices; author
      automation rules; accept or dismiss recommendations. &
      \textbf{Low–Medium} — can author rules (auto-action off by default) and control
      own unit's devices; cannot change the tariff or access another unit
      (\req{NFR-SEC-4}). \\
    \addlinespace
    Developer / Tester &
      Technical; prepares and reproduces demos. &
      Instantiate mock devices, exercise edge cases, force a device offline. &
      \textbf{Contained} — changes are limited to the mock world. \\
    \bottomrule
  \end{tabularx}
\end{table}

% -----------------------------------------------------------------------
\section{UI Requirements and the Seven Golden Rules}
\label{sec:ui-requirements}

\subsection{The seven golden rules applied to SHEO}

Seven widely-cited golden rules of interface design guide the evaluation.
Table~\ref{tab:golden-rules} lists each rule and its concrete realisation in SHEO.

\begin{longtable}{@{} L{3.0cm} p{9.2cm} @{}}
  \caption{The seven UI golden rules and their realisation in SHEO.}
  \label{tab:golden-rules} \\
  \toprule
  \textbf{Golden rule} & \textbf{Realisation in SHEO} \\
  \midrule
  \endfirsthead
  \multicolumn{2}{l}{\footnotesize\itshape (Table~\ref{tab:golden-rules} continued)} \\
  \toprule
  \textbf{Golden rule} & \textbf{Realisation in SHEO} \\
  \midrule
  \endhead
  \bottomrule
  \endfoot
  Easy to learn &
    The left navigation exposes every primary task from a single persistent menu;
    the dashboard lands on the most important information (live power and cycle savings)
    without any configuration.  A first-time Resident can operate a device in their unit
    and see live telemetry within a few minutes. \\
  \addlinespace
  Familiar terms &
    Domain language mirrors what residents already know: ``kWh today'', ``estimated
    bill'', and ``savings so far this cycle'' are expressed in the user's billing
    vocabulary rather than in software abstractions.  VND is used throughout for
    monetary figures rather than a generic currency symbol. \\
  \addlinespace
  Consistency &
    The \textbf{capability-driven control widget} is the central consistency benefit.
    Because one generic component renders every device type from the capability schema
    (\req{REQ-4.2.1}), the toggle, range slider, and segmented-enum interactions behave
    identically regardless of whether the user is controlling a plug, a bulb, a fan, or
    an air-conditioner.  There is no per-device-type UI code to learn or to diverge from
    each other. \\
  \addlinespace
  Minimal surprise &
    All destructive and automated actions are gated: auto-action on automation rules is
    \emph{off by default} (\req{BR-3}); the rule editor shows the VND estimate and any
    conflict warning \emph{before} the Save button is enabled; the Settings screen is
    read-only for non-Administrators. \\
  \addlinespace
  Recoverable &
    Two explicit recovery mechanisms are provided.  First, automation rules require the
    user to opt in to auto-apply; absent that opt-in the system notifies but takes no
    action, leaving the user in control.  Second, any rule firing that \emph{is}
    auto-applied can be undone within a two-minute window (\req{REQ-4.3.6}).
    Out-of-range or offline commands are rejected cleanly without changing device state. \\
  \addlinespace
  Feedback and help &
    Every command returns an explicit three-valued outcome: \textsc{success},
    \textsc{rejected}, or \textsc{timeout}.  The rule editor provides live validation and
    conflict warnings as the user types.  Offline devices are badged on the dashboard.
    The recommendations panel shows a rationale and a data-window label so users
    understand why each suggestion was generated.  WebSocket push ensures the dashboard
    reflects reality within five seconds (\req{NFR-PER-1}). \\
  \addlinespace
  Diverse users &
    The interface uses the Okabe--Ito colour-blind-safe palette throughout; the
    dashboard prominent headline-metric (large-numeral) pattern renders the key saving
    figure large enough to read at a glance; domain wording mirrors what residents already
    know; and the role-based access model ensures each actor sees only the controls
    relevant to their role, reducing cognitive load. \\
\end{longtable}

\subsection{Measurable usability targets}

In addition to the qualitative golden-rule analysis, the following measurable usability
targets were established from the context-of-use analysis and serve as the acceptance
criteria for the task-based evaluation in Section~\ref{sec:ui-evaluation}:

\begin{itemize}
  \item A user shall be able to send a control command to a device in $\le 3$ clicks
        from the dashboard (\req{REQ-4.2.x}, \req{NFR-PER-2}).
  \item UI feedback on every command outcome shall appear within the two-second
        budget (\req{NFR-PER-2}).
  \item A user shall be able to review a recommendation's rationale and VND figure
        and act on it in $\le 2$ clicks.
  \item The rule editor shall present the VND estimate and any conflict warning before
        the Save button becomes active (\req{REQ-4.5.3}).
  \item A first-time Resident shall be able to operate a device in their unit and see live
        telemetry within a few minutes (\req{REQ-4.2.5}, \req{REQ-4.1.1}).
\end{itemize}

\subsection{Colour and accessibility}

The Okabe--Ito palette is applied to all charts and status indicators.  The seven
Okabe--Ito colours are distinguishable for the three most common types of colour-vision
deficiency (deuteranopia, protanopia, and tritanopia), satisfying the ``diverse users''
golden rule at the perceptual level.

% -----------------------------------------------------------------------
\section{Information Architecture and Navigation}
\label{sec:ui-ia}

\subsection{Sitemap}

The application has a flat information architecture accessed through a
persistent left navigation bar, with a device-detail sub-view opened from the devices
screen.  The navigation is \textbf{role-aware}: a Resident or Developer sees the
unit-scoped Dashboard, Devices, Rules, Recommendations, and Reports screens, whereas the
Administrator sees a role-specific \textbf{Building} page (the units roster with totals
and a read-only drill-in) in place of the single-unit Dashboard.  After authentication,
every primary task for the actor's role is reachable in one click from the
navigation, and every primary task other than Settings is reachable within one additional
click from there.  The flat hierarchy minimises navigation depth, directly supporting the
``easy to learn'' golden rule.  Figure~\ref{fig:sitemap} shows the complete sitemap.

\fig{fig15-sitemap.png}{Sitemap of the SHEO single-page application.  After sign-in the shell exposes six top-level screens through the left navigation, plus a device-detail sub-view opened from the devices screen.  The rules screen opens a rule-editor overlay; the settings screen is read-only for non-Administrators (role-based access model, \req{BR-1}).  Every primary task is at most one click from the dashboard.}{fig:sitemap}

The Settings screen deserves special note: the building-wide tariff configuration and
unit-onboarding panels are visible to all authenticated users but are editable only by the
Administrator (\req{BR-1}).  Non-Administrators see the same layout in a read-only state
rather than a different screen, preserving layout consistency (``minimal surprise'').

\subsection{User flow}

The sitemap describes the static screen hierarchy; the user flow in
Figure~\ref{fig:userflow} describes the dynamic paths a user takes to accomplish the
three core journeys: controlling a device, acting on a recommendation, and authoring an
automation rule.  The diagram makes the key decision points explicit: the capability
fetch that precedes rendering, the three-valued command outcome, the accept-or-dismiss
branch on recommendations, and the live validation gate on the rule editor.

\fig{fig16-user-flow.png}{User flow for the three core journeys.  After sign-in the user lands on the Dashboard and chooses a journey: controlling a device (capability fetch $\rightarrow$ validate $\rightarrow$ \textsc{success}/\textsc{rejected}/\textsc{timeout}), acting on a recommendation (Accept $\Rightarrow$ becomes a rule; Dismiss $\Rightarrow$ 30-day suppression), or authoring a rule (live validation + conflict check + VND estimate before Save is enabled).}{fig:userflow}

\subsection{Three core journeys}

\begin{description}[style=nextline]
  \item[Journey 1 — Monitor and control a device.]
    From the Dashboard the user selects a device tile, the front-end requests the
    device's capability description from the server, and the control widget is rendered
    generically.  The user adjusts a control and the command flows through the
    command-application pipeline, which returns \textsc{success},
    \textsc{rejected}, or \textsc{timeout} with a human-readable message.

  \item[Journey 2 — Review and act on a recommendation.]
    The Recommendations screen presents up to five ranked suggestions, each showing a
    human-readable \emph{when\ldots then\ldots} rule expression, a rationale, the data window
    over which the habit was mined, and the estimated monthly saving in VND\@.
    Accepting a recommendation promotes it to a normal automation rule; dismissing it
    suppresses the same signature for 30 days (\req{REQ-4.4.5}).

  \item[Journey 3 — Author an automation rule.]
    The rule editor validates the rule in real time as the user constructs the
    \emph{when\ldots then\ldots until\ldots} expression.  Conflicting enabled rules
    are reported before the user can save.  The VND saving estimate is fetched from the
    server and displayed in the editor; the Save button is not enabled until a
    syntactically valid, conflict-free rule with a visible estimate is present.
\end{description}

% -----------------------------------------------------------------------
\section{Screen Design (Wireframes)}
\label{sec:ui-screens}

Six application screens were designed and prototyped as wireframes before implementation.
The wireframes are presented in Figures~\ref{fig:login}--\ref{fig:settings}.  All six
screens follow the same visual language: a left navigation rail, the page title and
primary action button in the top bar, and content laid out with a consistent spacing
grid.  The Okabe--Ito palette is applied to all charts and status indicators.

\medskip
A defining design constraint applies to the device-control screens: \textbf{every
control widget is rendered from the capability schema} (\req{REQ-4.2.1}).  The front-end
never branches on device type.  When the user navigates to a device's detail screen, the
front-end requests the device's capability description from the server and receives a
structured schema listing the device's controls — each classified as a toggle, a bounded
range, an enumeration, or a read-only readout — and the generic control component
renders the appropriate widget for each entry.
Adding a new device type to the system requires only a new capability schema entry —
no new UI code.

\figw{fig09e-login.png}{0.9}{Login screen.  Email and password sign-in establishes a session token held by the front-end for the duration of the session.  Four demo accounts (all with password \texttt{demo1234}) are accessible via pre-filled quick-login buttons, supporting the task-based evaluation: \texttt{admin@demo.com} (Administrator --- building owner with no unit, providing the building overview, tariff, and onboarding, and unable to operate devices), \texttt{resident@demo.com} (Resident of Unit~101, controlling that unit's devices), \texttt{resident2@demo.com} (Resident of Unit~102), and \texttt{dev@demo.com} (Developer/Tester on maintenance Unit~100).}{fig:login}

\figw{fig09a-dashboard.png}{0.9}{Dashboard.  Prominent headline-metric (large-numeral) displays show the unit's total power (W), kWh consumed today and in the current billing cycle, estimated bill in VND, and \textbf{savings so far this cycle in VND}.  A live WebSocket consumption chart, per-device tiles with offline badges, and a top-consumers panel complete the view (\req{REQ-4.1.1}--\req{REQ-4.1.5}, \req{REQ-4.5.4}).}{fig:dash}

\figw{fig09b-devices.png}{0.9}{Device detail screen.  Controls (on/off toggle, target temperature range, fan-speed enumeration) are generated entirely from the structured capability schema returned by the server — there is no per-device-type UI code.  Safety-sensitive controls carry a \emph{safety} warning pill (\req{NFR-SAF-3}); the fridge is safety-critical and not switchable, and the sensor is read-only.  The Resident operates only their own unit's operable devices, while the Developer/Tester can add or remove mock devices on the maintenance unit (\req{BR-1}).}{fig:devices}

\figw{fig09c-rules.png}{0.9}{Rule editor.  A \emph{when\ldots then\ldots until} rule form with live validation, conflict detection, and the VND saving estimate displayed before Save is enabled (\req{REQ-4.3.x}, \req{REQ-4.5.3}).  Auto-apply is off by default (\req{BR-3}); the execution history shows past firings with a two-minute undo link (\req{REQ-4.3.6}).}{fig:rules}

\figw{fig09d-recommendations.png}{0.9}{Recommendations screen.  Up to five ranked, human-readable suggestions are shown, each with the \emph{when\ldots then} rule expression, a rationale, the data window, and the estimated monthly saving in VND\@.  \emph{Accept} converts the recommendation into a normal rule; \emph{Dismiss} suppresses it for 30 days (\req{REQ-4.4.x}).}{fig:recs}

\figw{fig09f-settings.png}{0.9}{Settings screen.  The building-wide tariff configuration (flat, tiered EVN, or time-of-use; amounts in VND) and the unit-onboarding panel are editable only by the Administrator; non-Administrators see this screen in a read-only state (\req{BR-1}, \req{BR-5}, \req{NFR-SEC-4}).}{fig:settings}

\medskip
The wireframes above are deliberately \emph{low-fidelity} and monochrome: they fix
layout and information hierarchy before any visual styling is chosen.
Figure~\ref{fig:hifi} shows the corresponding \emph{high-fidelity} screens as actually
implemented --- in colour, on live seeded data --- completing the low- to
high-fidelity progression and evidencing that the design was carried through to a
working product.

\begin{figure}[htbp]
\centering
\begin{subfigure}{0.49\linewidth}\centering
  \includegraphics[width=\linewidth]{shot-dashboard.png}
  \caption{Dashboard}\label{fig:hifi-dash}
\end{subfigure}\hfill
\begin{subfigure}{0.49\linewidth}\centering
  \includegraphics[width=\linewidth]{shot-devices.png}
  \caption{Devices}\label{fig:hifi-devices}
\end{subfigure}

\medskip
\begin{subfigure}{0.49\linewidth}\centering
  \includegraphics[width=\linewidth]{shot-rules.png}
  \caption{Rules \& automation}\label{fig:hifi-rules}
\end{subfigure}\hfill
\begin{subfigure}{0.49\linewidth}\centering
  \includegraphics[width=\linewidth]{shot-recommendations.png}
  \caption{Recommendations}\label{fig:hifi-recs}
\end{subfigure}
\caption{High-fidelity realisation of four principal screens, captured from the running
application on seeded data. These implement the low-fidelity wireframes of
Figures~\ref{fig:dash}, \ref{fig:devices}, \ref{fig:rules} and~\ref{fig:recs}
respectively (\req{REQ-4.1}, \req{REQ-4.3}, \req{REQ-4.4}, \req{REQ-4.5}).}
\label{fig:hifi}
\end{figure}

% -----------------------------------------------------------------------
\section{Usability Evaluation}
\label{sec:ui-evaluation}

Following the evaluation step of the UI design process, the interface is assessed in two
passes: a formative heuristic walkthrough applied during development, and a summative
task-based test applied at the milestone demonstration.

\subsection{Heuristic walkthrough}

Each screen is examined against the seven golden rules from
Section~\ref{sec:ui-requirements}.  Rather than checking every screen independently, the
walkthrough focuses on the \emph{single interaction pattern that spans all device types}
— the capability-driven control widget — because validating one pattern validates the
interaction on every current and future device type simultaneously.

Two concrete design defects surfaced during this pass and were corrected before the
implementation milestone:

\begin{enumerate}
  \item \textbf{VND estimate placement.}  An early iteration placed the saving estimate
        below the Save button, making it easy to miss.  The corrected design disables
        the Save button until the estimate has been fetched and displayed, forcing the
        user to see the figure before committing.
  \item \textbf{Safety visibility.}  Safety-sensitive controls (such as the AC
        target-temperature control) were initially indistinguishable from ordinary controls.  The corrected
        design adds a prominent \emph{safety} warning pill adjacent to any control that
        the capability schema identifies as safety-sensitive (\req{NFR-SAF-3}).
\end{enumerate}

The heuristic pass also confirmed that auto-action defaulting to off (\req{BR-3}) and
the two-minute undo window (\req{REQ-4.3.6}) together satisfy the ``recoverable'' golden
rule for all three core journeys.

\subsection{Task-based usability test}

At the milestone demonstration, evaluators perform the four tasks in
Table~\ref{tab:usability-tasks} using the four demo accounts, spanning the
Administrator, Resident, and Developer/Tester roles.  For each task, successful completion is defined by a measurable
criterion; any task that fails its criterion is treated as a defect and triggers a design
revision before release.

\begin{table}[htbp]
  \caption{Task-based usability test: tasks and measurable success criteria.}
  \label{tab:usability-tasks}
  \begin{tabularx}{\linewidth}{@{} L{3.2cm} Y L{2.8cm} @{}}
    \toprule
    \textbf{Task (journey)} & \textbf{Measurable success criterion} & \textbf{Tied to} \\
    \midrule
    Control a device &
      Reach a device's controls and send a command in $\le 3$ clicks from the dashboard;
      the outcome (\textsc{success} / \textsc{rejected} / \textsc{timeout}) is shown
      within the two-second UI-feedback budget. &
      \req{REQ-4.2.x},\newline \req{NFR-PER-2} \\
    \addlinespace
    Review and act on a recommendation &
      Read the \emph{when\ldots then} rule expression, rationale, and VND figure, then
      Accept or Dismiss in $\le 2$ clicks; the VND figure is visible without scrolling. &
      \req{REQ-4.4.x} \\
    \addlinespace
    Author a rule &
      Build a valid rule and see its VND estimate \emph{before} Save is enabled; a
      conflicting enabled rule is reported inline, not silently overwritten. &
      \req{REQ-4.3.x},\newline \req{REQ-4.5.3} \\
    \addlinespace
    First-run onboarding &
      A first-time Resident operates a device in their unit and observes live telemetry on
      the dashboard within a few minutes of signing in to the unit. &
      \req{REQ-4.2.5},\newline \req{REQ-4.1.1} \\
    \bottomrule
  \end{tabularx}
\end{table}

Error prevention and recoverability are tested explicitly as part of each task: an
out-of-range or offline command must surface a clear rejection message and leave device
state unchanged; the auto-action toggle must default to off; and any auto-applied rule
firing must be undoable for two minutes.  Evaluation results are written up in a
brief usability report; any failing task loops back into the design (re-layout, clearer
labels, or additional feedback wording) before the release build is produced.
